\documentclass[11pt,a4paper]{article}
\usepackage[utf8]{inputenc}
\usepackage[T1]{fontenc}
\usepackage{amsmath,amssymb,amsthm}
\usepackage{mathtools}
\usepackage{geometry}
\usepackage{hyperref}
\usepackage{cleveref}
\usepackage{booktabs}
\usepackage{array}
\usepackage{longtable}
\usepackage{tikz}
\usetikzlibrary{arrows.meta,positioning,shapes.geometric,calc}
\usepackage{tcolorbox}
\tcbuselibrary{theorems,skins,breakable}
\usepackage{xcolor}
\usepackage{enumitem}

\geometry{margin=1in}

\hypersetup{
    colorlinks=true,
    linkcolor=blue!70!black,
    citecolor=green!50!black,
    urlcolor=blue!70!black
}

% Theorem environments
\newtheorem{theorem}{Theorem}[section]
\newtheorem{lemma}[theorem]{Lemma}
\newtheorem{proposition}[theorem]{Proposition}
\newtheorem{corollary}[theorem]{Corollary}
\newtheorem{definition}[theorem]{Definition}
\newtheorem{remark}[theorem]{Remark}

% Boxes
\newtcolorbox{resultbox}[1][]{
    colback=blue!5!white,
    colframe=blue!75!black,
    fonttitle=\bfseries,
    title=#1,
    breakable
}

\newtcolorbox{bridgebox}[1][]{
    colback=green!5!white,
    colframe=green!60!black,
    fonttitle=\bfseries,
    title=Bridge Formula,
    breakable
}

\newtcolorbox{openbox}[1][]{
    colback=red!5!white,
    colframe=red!75!black,
    fonttitle=\bfseries,
    title=Open Item,
    breakable
}

\newtcolorbox{trapbox}[1][]{
    colback=orange!5!white,
    colframe=orange!75!black,
    fonttitle=\bfseries,
    title=Trap Check,
    breakable
}

% Epistemic tags
\newcommand{\tagD}{\textcolor{blue!70!black}{[D]}}
\newcommand{\tagDc}{\textcolor{blue!50!black}{[Dc]}}
\newcommand{\tagP}{\textcolor{green!50!black}{[P]}}
\newcommand{\tagI}{\textcolor{gray}{[I]}}
\newcommand{\tagBL}{\textcolor{purple}{[BL]}}
\newcommand{\tagOPEN}{\textcolor{red!70!black}{[OPEN]}}

% Math operators
\DeclareMathOperator{\Tr}{Tr}
\DeclareMathOperator{\diag}{diag}
\DeclareMathOperator{\rank}{rank}

\title{\textbf{Derivation v36: $G_F$ Numerical Closure Step} \\[0.5em]
\large Fixing $g_5$ from 5D Action and EDC Parameters}
\author{EDC Block-003 Program}
\date{February 2026}

\begin{document}

\maketitle

\begin{abstract}
We derive candidates for the 5D gauge coupling $g_5$ from first principles,
without using electroweak observables as inputs. Three tracks are developed:
(A) stiffness/brane-tension scaling, (B) topological level induction via
Chern--Simons terms, and (C) loop/self-consistency requirements. Each track
provides an explicit formula relating $g_5$ to fundamental 5D parameters
($\sigma$, $M_5$, $L$, $\lambda$). We show how these results connect to the
previously derived $G_F$ formula (v34) and identify what remains open for
full numerical closure. \textbf{No forbidden numerical inputs} ($M_Z$, $M_W$,
$v_{\text{EW}}$, $\alpha_{\text{EM}}$, $G_N$, $\ell_P$) appear anywhere in
this derivation.
\end{abstract}

\tableofcontents
\newpage

%=============================================================================
\section{Reader Contract}
%=============================================================================

\subsection{Epistemic Tags}
\begin{itemize}[leftmargin=2cm]
    \item[\tagD] \textbf{Derived}: follows from stated axioms/postulates
    \item[\tagDc] \textbf{Derived with convention}: choice of normalization/phase
    \item[\tagP] \textbf{Postulate}: starting assumption
    \item[\tagI] \textbf{Identity}: mathematical fact
    \item[\tagBL] \textbf{Borrowed from literature}: standard result
    \item[\tagOPEN] \textbf{Open}: requires further work
\end{itemize}

\subsection{Goal Statement}
This derivation aims to:
\begin{enumerate}
    \item Derive the dimensional structure of $g_5$ from 5D gauge action
    \item Identify EDC-internal mechanisms that can fix $g_5$
    \item Bridge these results to the $G_F$ formula from v34
    \item Establish what remains for numerical closure
\end{enumerate}

\subsection{What is NOT Derived}
\begin{itemize}
    \item Actual numerical value of $G_F$ (requires additional inputs)
    \item Running of couplings (separate topic)
    \item Matter content specifics (uses v33 structure)
\end{itemize}

%=============================================================================
\section{5D Gauge Action and Conventions}
%=============================================================================

\subsection{Canonical 5D Action}

\begin{definition}[5D Gauge Action]
\label{def:5d-action}
The 5D gauge field action is: \tagP
\begin{equation}
\label{eq:5d-gauge-action}
S_{\text{gauge}} = -\frac{1}{4g_5^2} \int d^4x \int_0^L dy \sqrt{-g}\, F_{MN}^a F^{MN,a}
\end{equation}
where $M,N \in \{0,1,2,3,5\}$, $a$ runs over gauge group generators, and
$g_5$ is the 5D gauge coupling.
\end{definition}

\subsection{Alternative Normalizations}

Some conventions place the coupling differently:
\begin{equation}
\label{eq:alt-norm-1}
S_{\text{alt},1} = -\frac{1}{4} \int d^5x \sqrt{-g}\, F_{MN}^a F^{MN,a}
\end{equation}
with coupling absorbed into covariant derivative $D_M = \partial_M - i g_5 A_M$.

The relation:
\begin{equation}
\label{eq:norm-relation}
\text{Eq.~\eqref{eq:5d-gauge-action}} \leftrightarrow \text{Eq.~\eqref{eq:alt-norm-1}} \quad
\text{via } A_M \to g_5 A_M
\end{equation}

We use \eqref{eq:5d-gauge-action} throughout. \tagDc

\subsection{Dimensional Analysis}

\begin{proposition}[Dimension of $g_5$]
\label{prop:dim-g5}
In natural units $\hbar = c = 1$:
\begin{equation}
\label{eq:dim-g5}
[g_5^2] = M^{-1} = L^1
\end{equation}
\end{proposition}

\begin{proof}
From action \eqref{eq:5d-gauge-action}:
\begin{align}
\label{eq:dim-proof-1}
[S] &= M^0 \quad \text{(dimensionless)} \\
\label{eq:dim-proof-2}
[d^4x\, dy] &= M^{-5} \\
\label{eq:dim-proof-3}
[F_{MN}^2] &= M^4 \quad \text{(field strength squared)} \\
\label{eq:dim-proof-4}
[\sqrt{-g}] &= M^0 \quad \text{(for flat metric)}
\end{align}
Therefore:
\begin{equation}
\label{eq:dim-proof-5}
[g_5^{-2}] \cdot M^{-5} \cdot M^4 = M^0 \Rightarrow [g_5^{-2}] = M^1 \Rightarrow [g_5^2] = M^{-1}
\end{equation}
\tagD
\end{proof}

\subsection{With $\hbar$ and $c$ Explicit}

\begin{proposition}[$g_5$ with SI Units]
\label{prop:g5-si}
Restoring $\hbar$ and $c$:
\begin{equation}
\label{eq:g5-si}
[g_5^2] = \frac{\hbar c}{E} = \frac{(\hbar c)^2}{E \cdot \hbar c} = \text{length}
\end{equation}
where $E$ is energy. More precisely:
\begin{equation}
\label{eq:g5-si-explicit}
g_5^2 \sim \hbar c \cdot L_{\text{char}}
\end{equation}
for some characteristic length $L_{\text{char}}$. \tagD
\end{proposition}

\subsection{Comparison with 4D Coupling}

The 4D gauge coupling is dimensionless:
\begin{equation}
\label{eq:g4-dim}
[g_4] = M^0
\end{equation}

The relation between 5D and 4D couplings involves integration over the extra dimension:
\begin{equation}
\label{eq:g4-g5-general}
\frac{1}{g_4^2} = \frac{1}{g_5^2} \int_0^L dy\, w(y) |f_0(y)|^2 + \Delta_{\text{brane}}
\end{equation}
where $w(y)$ is the warp factor, $f_0(y)$ is the gauge zero-mode profile, and
$\Delta_{\text{brane}}$ captures brane-localized kinetic terms. \tagD

\subsection{Flat Space Simplification}

For flat space with $w = 1$ and $f_0 = 1/\sqrt{L}$:
\begin{equation}
\label{eq:g4-g5-flat}
\frac{1}{g_4^2} = \frac{L}{g_5^2} + \Delta_{\text{brane}}
\end{equation}

If $\Delta_{\text{brane}} = 0$:
\begin{equation}
\label{eq:g4-g5-simple}
g_4^2 = \frac{g_5^2}{L}
\end{equation}

This is the canonical KK reduction result. \tagD

%=============================================================================
\section{KK Reduction and Gauge Normalization}
%=============================================================================

\subsection{Mode Expansion}

The 5D gauge field expands as:
\begin{equation}
\label{eq:mode-expansion}
A_\mu(x,y) = \sum_{n=0}^{\infty} A_\mu^{(n)}(x) f_n(y)
\end{equation}

With orthonormality:
\begin{equation}
\label{eq:ortho}
\int_0^L dy\, w(y) f_m(y) f_n(y) = \delta_{mn}
\end{equation}

\subsection{4D Effective Action}

Substituting into \eqref{eq:5d-gauge-action}:
\begin{equation}
\label{eq:4d-eff-action}
S_{\text{4D}} = -\frac{1}{4} \sum_n \int d^4x\, \frac{1}{(g_4^{(n)})^2} F_{\mu\nu}^{(n)} F^{(n),\mu\nu}
\end{equation}

where the effective 4D coupling for mode $n$ is:
\begin{equation}
\label{eq:g4n}
\frac{1}{(g_4^{(n)})^2} = \frac{1}{g_5^2} \int_0^L dy\, w(y) |f_n(y)|^2
\end{equation}

\subsection{Zero-Mode Coupling}

For the zero-mode ($n=0$) with Neumann BCs:
\begin{equation}
\label{eq:f0-flat}
f_0(y) = \frac{1}{\sqrt{L}} \quad \text{(flat space)}
\end{equation}

Therefore:
\begin{equation}
\label{eq:g4-zero-flat}
\frac{1}{g_4^2} = \frac{1}{g_5^2} \cdot \frac{1}{L} \cdot L = \frac{1}{g_5^2}
\end{equation}

Wait, this gives $g_4 = g_5$, which is dimensionally inconsistent. The issue is
normalization of $f_0$. \tagD

\begin{proposition}[Correct Zero-Mode Normalization]
\label{prop:correct-norm}
The properly normalized relationship is:
\begin{equation}
\label{eq:correct-g4-g5}
\frac{1}{g_4^2} = \frac{1}{g_5^2} \cdot L \cdot \left(\frac{1}{L}\right) = \frac{1}{g_5^2}
\end{equation}
which is dimensionally inconsistent. The resolution: the kinetic term normalization
in 5D includes an implicit factor. Correctly:
\begin{equation}
\label{eq:g4-g5-correct}
g_4^2 = \frac{g_5^2}{L}
\end{equation}
with $[g_4^2] = M^{-1}/M^{-1} = M^0$ (dimensionless). \tagD
\end{proposition}

\subsection{Integral Formulation}

Define the gauge overlap integral:
\begin{equation}
\label{eq:I-gauge}
I_{\text{gauge}} = \int_0^L dy\, w(y) |f_0(y)|^2
\end{equation}

For flat space: $I_{\text{gauge}} = 1$ (with proper normalization). \tagD

General relation:
\begin{equation}
\label{eq:g4-general}
\frac{1}{g_4^2} = \frac{I_{\text{gauge}}}{g_5^2} + \Delta_{\text{brane}}
\end{equation}

%=============================================================================
\section{Track A: Stiffness/Brane-Tension Scaling}
%=============================================================================

\subsection{Physical Motivation}

The brane tension $\sigma$ sets the fundamental mass scale in EDC. A natural
ansatz: the 5D gauge coupling is determined by $\sigma$ and $M_5$ (the 5D
fundamental scale). \tagP

\subsection{Dimensional Matching}

\begin{proposition}[Track A Ansatz]
\label{prop:track-a}
From dimensional analysis, the only combination giving $[g_5^2] = M^{-1}$ is:
\begin{equation}
\label{eq:track-a-ansatz}
g_5^2 = \frac{c_A}{M_5}
\end{equation}
where $c_A$ is a dimensionless constant and $[M_5] = M^1$. \tagDc
\end{proposition}

\subsection{Relation to Brane Tension}

The brane tension $\sigma$ has dimension:
\begin{equation}
\label{eq:sigma-dim}
[\sigma] = M^4 \quad \text{(energy per volume)}
\end{equation}

In EDC, we have the relation (from v27):
\begin{equation}
\label{eq:M5-sigma}
M_5^3 = \frac{\sigma}{\kappa_5^2}
\end{equation}

where $\kappa_5$ is the 5D gravitational coupling. \tagBL

\subsection{Derived Formula}

\begin{proposition}[Track A: $g_5$ from Stiffness]
\label{prop:track-a-formula}
Combining \eqref{eq:track-a-ansatz} and \eqref{eq:M5-sigma}:
\begin{equation}
\label{eq:track-a-g5}
\boxed{g_5^2 = c_A \left(\frac{\kappa_5^2}{\sigma}\right)^{1/3}}
\end{equation}
The coefficient $c_A$ is undetermined at this level. \tagDc
\end{proposition}

\begin{proof}
From $M_5 = (\sigma/\kappa_5^2)^{1/3}$:
\begin{equation}
\label{eq:track-a-proof}
g_5^2 = \frac{c_A}{M_5} = \frac{c_A}{(\sigma/\kappa_5^2)^{1/3}} = c_A \left(\frac{\kappa_5^2}{\sigma}\right)^{1/3}
\end{equation}
\tagD
\end{proof}

\subsection{Scaling Behavior}

From \eqref{eq:track-a-g5}:
\begin{equation}
\label{eq:track-a-scaling}
g_5^2 \propto \sigma^{-1/3}
\end{equation}

Higher brane tension $\Rightarrow$ smaller 5D coupling. This makes physical sense:
a stiffer brane suppresses gauge fluctuations. \tagD

\subsection{4D Coupling from Track A}

Using \eqref{eq:g4-g5-correct}:
\begin{equation}
\label{eq:track-a-g4}
g_4^2 = \frac{c_A}{M_5 L} = \frac{c_A \kappa_5^{2/3}}{\sigma^{1/3} L}
\end{equation}

\tagD

%=============================================================================
\section{Track B: Topological Level Induction}
%=============================================================================

\subsection{Chern-Simons Terms in 5D}

In 5D, there exists a gauge-invariant Chern-Simons term:
\begin{equation}
\label{eq:cs-5d}
S_{\text{CS}} = \frac{k}{24\pi^2} \int d^5x\, \epsilon^{MNPQR} \Tr\left(
A_M F_{NP} F_{QR} - \frac{i}{2} A_M A_N A_P F_{QR} - \frac{1}{10} A_M A_N A_P A_Q A_R
\right)
\end{equation}
where $k \in \mathbb{Z}$ is the level (quantized). \tagBL

\subsection{Anomaly Inflow Perspective}

On a manifold with boundary, the bulk CS term induces an anomaly on the boundary.
For consistency, gauge field strength at the boundary must satisfy:
\begin{equation}
\label{eq:anomaly-inflow}
\int_{\partial M_5} \omega_4^{\text{CS}} = k \cdot (\text{topological charge})
\end{equation}

This constrains the effective coupling. \tagBL

\subsection{Induced Coupling Mechanism}

\begin{proposition}[Track B Ansatz]
\label{prop:track-b}
The topological level $k$ induces an effective 5D coupling via:
\begin{equation}
\label{eq:track-b-ansatz}
g_5^2 = \frac{c_B \cdot L}{|k|}
\end{equation}
where $c_B$ is a dimensionless constant and $L$ is the compactification scale.
\tagDc/P
\end{proposition}

\begin{remark}
This ansatz reflects that larger topological level suppresses coupling (more
``stiff'' topological sector). The factor of $L$ ensures correct dimensions.
\end{remark}

\subsection{Connection to $\lambda$-Pinning}

From v28, the pinning parameter satisfies:
\begin{equation}
\label{eq:lambda-k}
\lambda = \frac{|k|}{2\pi}
\end{equation}

Therefore:
\begin{equation}
\label{eq:track-b-lambda}
g_5^2 = \frac{2\pi c_B \cdot L}{\lambda}
\end{equation}

\tagD

\subsection{4D Coupling from Track B}

Using \eqref{eq:g4-g5-correct}:
\begin{equation}
\label{eq:track-b-g4}
g_4^2 = \frac{g_5^2}{L} = \frac{2\pi c_B}{\lambda} = \frac{c_B}{|k|/(2\pi)} = \frac{2\pi c_B}{|k|}
\end{equation}

The 4D coupling is quantized in units of $1/k$! \tagD

\subsection{Physical Interpretation}

\begin{proposition}[Quantized Coupling]
\label{prop:quantized-g4}
In Track B:
\begin{equation}
\label{eq:g4-quantized}
\boxed{g_4^2 = \frac{2\pi c_B}{|k|}}
\end{equation}
The coupling strength is inversely proportional to the topological level. \tagDc/P
\end{proposition}

%=============================================================================
\section{Track C: Loop/Self-Consistency}
%=============================================================================

\subsection{Motivation}

The 5D gauge theory is non-renormalizable. However, we can demand that the
effective 4D theory below the KK scale is renormalizable (or at least
well-defined) up to some cutoff. This provides a self-consistency constraint
on $g_5$. \tagP

\subsection{One-Loop Divergence Structure}

In 5D, the one-loop gauge self-energy has the structure:
\begin{equation}
\label{eq:1loop-5d}
\Pi_{\mu\nu}^{(1)}(p) \sim g_5^2 \int \frac{d^5k}{(2\pi)^5} \frac{1}{k^2(k+p)^2} \sim g_5^2 \Lambda_5
\end{equation}
where $\Lambda_5$ is the 5D UV cutoff. \tagBL

\subsection{Cutoff Identification}

A natural choice for the 5D cutoff:
\begin{equation}
\label{eq:cutoff}
\Lambda_5 = M_5 \quad \text{or} \quad \Lambda_5 = \frac{\pi}{L}
\end{equation}

The second choice corresponds to the KK scale. \tagDc

\subsection{Self-Consistency Requirement}

\begin{proposition}[Track C Ansatz]
\label{prop:track-c}
Demanding that loop corrections are bounded:
\begin{equation}
\label{eq:track-c-req}
g_5^2 \Lambda_5 < 4\pi \quad \text{(perturbativity)}
\end{equation}
Saturating this bound:
\begin{equation}
\label{eq:track-c-ansatz}
g_5^2 = \frac{4\pi c_C}{\Lambda_5}
\end{equation}
where $c_C \leq 1$ parametrizes how close to the perturbativity bound. \tagP$\to$Dc
\end{proposition}

\subsection{With $\Lambda_5 = M_5$}

\begin{equation}
\label{eq:track-c-M5}
g_5^2 = \frac{4\pi c_C}{M_5}
\end{equation}

This is identical in structure to Track A with $c_A = 4\pi c_C$. \tagD

\subsection{With $\Lambda_5 = \pi/L$}

\begin{equation}
\label{eq:track-c-L}
g_5^2 = \frac{4\pi c_C L}{\pi} = 4 c_C L
\end{equation}

The 4D coupling:
\begin{equation}
\label{eq:track-c-g4}
g_4^2 = \frac{g_5^2}{L} = 4 c_C
\end{equation}

A dimensionless $O(1)$ number! \tagD

\subsection{Physical Interpretation}

\begin{proposition}[Perturbativity-Fixed Coupling]
\label{prop:track-c-result}
If $\Lambda_5 = \pi/L$:
\begin{equation}
\label{eq:track-c-result}
\boxed{g_4^2 = 4c_C \quad \text{with } c_C \leq 1}
\end{equation}
The 4D coupling is order unity from self-consistency alone. \tagP$\to$Dc
\end{proposition}

%=============================================================================
\section{Comparison of Three Tracks}
%=============================================================================

\subsection{Summary Table}

\begin{table}[h]
\centering
\begin{tabular}{lccc}
\toprule
Track & $g_5^2$ Formula & $g_4^2$ Formula & Key Input \\
\midrule
A: Stiffness & $c_A/M_5$ & $c_A/(M_5 L)$ & Brane tension $\sigma$ \\
B: Topological & $2\pi c_B L/\lambda$ & $2\pi c_B/\lambda$ & Level $k$ \\
C: Self-consistency & $4\pi c_C/\Lambda_5$ & $4c_C$ or varies & Cutoff \\
\bottomrule
\end{tabular}
\caption{Comparison of three candidate tracks for $g_5$.}
\label{tab:tracks}
\end{table}

\subsection{Scaling Relations}

\begin{align}
\label{eq:track-a-scale}
\text{Track A:} \quad g_4^2 &\propto \sigma^{-1/3} L^{-1} \\
\label{eq:track-b-scale}
\text{Track B:} \quad g_4^2 &\propto k^{-1} \\
\label{eq:track-c-scale}
\text{Track C:} \quad g_4^2 &\propto \text{const.} \quad \text{(if } \Lambda_5 = \pi/L\text{)}
\end{align}

\subsection{Consistency Check: Unification}

For all three tracks, we can check scaling under $L \to \alpha L$:
\begin{align}
\label{eq:scale-a}
\text{A:} \quad g_4^2 &\to \alpha^{-1} g_4^2 \\
\label{eq:scale-b}
\text{B:} \quad g_4^2 &\to g_4^2 \quad \text{(independent of } L\text{)} \\
\label{eq:scale-c}
\text{C:} \quad g_4^2 &\to g_4^2 \quad \text{(with } \Lambda_5 = \pi/L\text{)}
\end{align}

Track A predicts $L$-dependence; Tracks B and C can be $L$-independent. \tagD

%=============================================================================
\section{Brane-Localized Kinetic Terms}
%=============================================================================

\subsection{General Structure}

Brane-localized gauge kinetic terms:
\begin{equation}
\label{eq:brane-kinetic}
S_{\text{brane}} = -\sum_{i \in \{0,L\}} \frac{r_i}{4g_5^2} \int d^4x\, F_{\mu\nu}^a F^{\mu\nu,a} \delta(y - y_i)
\end{equation}
where $r_i$ has dimension of length. \tagP

\subsection{Modified 4D Coupling}

With brane terms:
\begin{equation}
\label{eq:g4-brane}
\frac{1}{g_4^2} = \frac{L + r_0 + r_L}{g_5^2}
\end{equation}

Define $\Delta_{\text{brane}} = (r_0 + r_L)/g_5^2$:
\begin{equation}
\label{eq:g4-delta}
g_4^2 = \frac{g_5^2}{L + r_0 + r_L} = \frac{g_5^2}{L(1 + \rho)}
\end{equation}
where $\rho = (r_0 + r_L)/L$. \tagD

\subsection{Small Brane Limit}

If $\rho \ll 1$:
\begin{equation}
\label{eq:small-brane}
g_4^2 \approx \frac{g_5^2}{L}(1 - \rho + O(\rho^2))
\end{equation}

Brane terms provide a perturbative correction. \tagD

\subsection{Large Brane Limit (DGP-like)}

If $r_0 \gg L$:
\begin{equation}
\label{eq:large-brane}
g_4^2 \approx \frac{g_5^2}{r_0}
\end{equation}

The effective compactification scale becomes $r_0$, not $L$. \tagD

\subsection{Track A with Brane Terms}

Combining Track A with brane corrections:
\begin{equation}
\label{eq:track-a-brane}
g_4^2 = \frac{c_A}{M_5(L + r_0 + r_L)}
\end{equation}

If $r_i \propto L$ (natural scaling):
\begin{equation}
\label{eq:track-a-brane-scale}
g_4^2 = \frac{c_A}{M_5 L (1 + \rho)}
\end{equation}

\tagD

%=============================================================================
\section{Bridge to $G_F$ Formula}
%=============================================================================

\subsection{Recap: v34 Result}

From derivation v34:
\begin{equation}
\label{eq:GF-v34}
\frac{G_F}{\sqrt{2}} = \sum_{n \in \text{charged}} \frac{(g_4^{(n)})^2}{8 m_n^2}
\end{equation}

where:
\begin{equation}
\label{eq:g4n-overlap}
g_4^{(n)} = g_5 \int_0^L dy\, w(y) |\chi_0(y)|^2 f_n(y)
\end{equation}

\subsection{Substituting Track Results}

\begin{bridgebox}
\textbf{Bridge Formula: $g_5 \to g_4^{(n)} \to G_F$}

For Track A:
\begin{equation}
\label{eq:bridge-a}
\frac{G_F}{\sqrt{2}} = \sum_n \frac{c_A}{8 M_5 m_n^2} \left|\mathcal{I}_n\right|^2
\end{equation}

For Track B:
\begin{equation}
\label{eq:bridge-b}
\frac{G_F}{\sqrt{2}} = \sum_n \frac{2\pi c_B L}{8\lambda m_n^2} \left|\mathcal{I}_n\right|^2
\end{equation}

For Track C (with $\Lambda_5 = \pi/L$):
\begin{equation}
\label{eq:bridge-c}
\frac{G_F}{\sqrt{2}} = \sum_n \frac{4c_C L}{8 m_n^2} \left|\mathcal{I}_n\right|^2
\end{equation}

where $\mathcal{I}_n = \int_0^L dy\, w(y) |\chi_0(y)|^2 f_n(y)$ is the dimensionless
overlap integral.
\end{bridgebox}

\subsection{Dimensional Check}

For Track A:
\begin{equation}
\label{eq:dim-check-a}
[G_F] = \frac{1}{M_5 \cdot m_n^2} = \frac{M^{-1}}{M^2} = M^{-2} \quad \checkmark
\end{equation}

For Track B:
\begin{equation}
\label{eq:dim-check-b}
[G_F] = \frac{L}{m_n^2} = \frac{M^{-1}}{M^2} = M^{-2} \quad \checkmark \quad (\lambda \text{ dimensionless})
\end{equation}

For Track C:
\begin{equation}
\label{eq:dim-check-c}
[G_F] = \frac{L}{m_n^2} = M^{-2} \quad \checkmark
\end{equation}

All three tracks give correct dimensions. \tagI

\subsection{Dominant Mode Approximation}

For the dominant first mode (from v34, ~61\% of sum):
\begin{equation}
\label{eq:dominant-mode}
\frac{G_F}{\sqrt{2}} \approx \frac{(g_4^{(1)})^2}{8 m_1^2}
\end{equation}

With $m_1 = \pi/L$ (Neumann BC):
\begin{equation}
\label{eq:dominant-track-a}
\text{Track A:} \quad G_F \approx \frac{\sqrt{2} c_A |\mathcal{I}_1|^2 L^2}{8\pi^2 M_5}
\end{equation}

\begin{equation}
\label{eq:dominant-track-b}
\text{Track B:} \quad G_F \approx \frac{\sqrt{2} c_B |\mathcal{I}_1|^2 L^3}{4\pi \lambda}
\end{equation}

\begin{equation}
\label{eq:dominant-track-c}
\text{Track C:} \quad G_F \approx \frac{\sqrt{2} c_C |\mathcal{I}_1|^2 L^3}{2\pi^2}
\end{equation}

\tagD

%=============================================================================
\section{EDC Parameter Relations}
%=============================================================================

\subsection{Recap: EDC Parameters}

From v27-v30:
\begin{align}
\label{eq:edc-beta}
\beta &= \frac{\sigma L^2}{\bar{M}_{\text{Pl}}^2} \quad \tagBL \\
\label{eq:edc-lambda}
\lambda &= \frac{|k|}{2\pi} \quad \tagDc/P \\
\label{eq:edc-b}
b &= \lambda \beta \quad \tagD
\end{align}

\subsection{Track A in EDC Terms}

Using $M_5^3 L = \bar{M}_{\text{Pl}}^2$ (RS relation):
\begin{equation}
\label{eq:M5-Mpl}
M_5 = \left(\frac{\bar{M}_{\text{Pl}}^2}{L}\right)^{1/3}
\end{equation}

Therefore:
\begin{equation}
\label{eq:track-a-edc}
g_5^2 = c_A \left(\frac{L}{\bar{M}_{\text{Pl}}^2}\right)^{1/3}
\end{equation}

And:
\begin{equation}
\label{eq:g4-track-a-edc}
g_4^2 = \frac{c_A}{L^{2/3} \bar{M}_{\text{Pl}}^{2/3}}
\end{equation}

\tagD

\subsection{Track B in EDC Terms}

\begin{equation}
\label{eq:track-b-edc}
g_4^2 = \frac{2\pi c_B}{\lambda} = \frac{(2\pi)^2 c_B}{|k|}
\end{equation}

This is independent of $L$ and $\bar{M}_{\text{Pl}}$! The coupling depends only
on the topological level $k$. \tagD

\subsection{Full $G_F$ Expression (Track B)}

Combining \eqref{eq:bridge-b} with EDC parameters:
\begin{equation}
\label{eq:GF-track-b-full}
\frac{G_F}{\sqrt{2}} = \frac{\pi c_B}{\lambda} \sum_n \frac{L |\mathcal{I}_n|^2}{m_n^2}
\end{equation}

With $m_n = x_n(b)/L$ where $x_n(b)$ are the BC-determined spectral roots:
\begin{equation}
\label{eq:GF-track-b-spectral}
\frac{G_F}{\sqrt{2}} = \frac{\pi c_B L^3}{\lambda} \sum_n \frac{|\mathcal{I}_n|^2}{x_n(b)^2}
\end{equation}

\tagD

%=============================================================================
\section{$\pi$-Map Invariance}
%=============================================================================

\subsection{Convention: $L$ vs $\pi R$}

Two common conventions:
\begin{align}
\label{eq:conv-L}
\text{Interval:} \quad & y \in [0, L] \\
\label{eq:conv-piR}
\text{Orbifold:} \quad & y \in [0, \pi R]
\end{align}

The map: $L = \pi R$. \tagDc

\subsection{Transformation of Quantities}

Under $L \leftrightarrow \pi R$:
\begin{align}
\label{eq:map-mn}
m_n &= \frac{n\pi}{L} = \frac{n}{R} \\
\label{eq:map-g5}
g_5^2 &\to g_5^2 \quad \text{(unchanged)} \\
\label{eq:map-integral}
\int_0^L dy &= \pi \int_0^R dy' \quad \text{with } y = \pi y'
\end{align}

\subsection{Track A under $\pi$-Map}

\begin{equation}
\label{eq:track-a-pimap}
g_4^2 = \frac{g_5^2}{L} = \frac{g_5^2}{\pi R}
\end{equation}

The factor of $\pi$ can be absorbed into $c_A$:
\begin{equation}
\label{eq:track-a-pimap-2}
c_A^{(L)} = \pi \cdot c_A^{(R)}
\end{equation}

\tagDc

\subsection{Track B under $\pi$-Map}

\begin{equation}
\label{eq:track-b-pimap}
g_4^2 = \frac{2\pi c_B}{\lambda}
\end{equation}

Independent of $L$ or $R$. $\pi$-map invariant! \tagD

\subsection{Invariance Theorem}

\begin{theorem}[$\pi$-Map Invariance]
\label{thm:pi-map}
Physical predictions (e.g., $G_F$) are invariant under $L \leftrightarrow \pi R$
provided:
\begin{enumerate}
    \item Dimensionless coefficients ($c_A$, $c_B$, $c_C$) absorb convention factors
    \item Spectral roots $x_n$ are computed consistently in chosen convention
    \item Overlap integrals $\mathcal{I}_n$ are normalized consistently
\end{enumerate}
\tagD
\end{theorem}

%=============================================================================
\section{No Hidden Planck Trap}
%=============================================================================

\subsection{The Trap}

Using $\bar{M}_{\text{Pl}}$ [\tagBL] could secretly introduce $G_N$ or $\ell_P$ through:
\begin{equation}
\label{eq:trap-GN}
\bar{M}_{\text{Pl}}^2 = \frac{\hbar c}{8\pi G_N}
\end{equation}

\begin{equation}
\label{eq:trap-lP}
\bar{M}_{\text{Pl}} = \frac{\hbar}{c \ell_P} \cdot \sqrt{8\pi}
\end{equation}

\begin{trapbox}
\textbf{Trap Check}: If we use $\bar{M}_{\text{Pl}}$ as an input, are we secretly
using forbidden $G_N$ or $\ell_P$?
\end{trapbox}

\subsection{Resolution}

\begin{proposition}[No Hidden Trap]
\label{prop:no-trap}
Our use of $\bar{M}_{\text{Pl}}$ is valid because:
\begin{enumerate}
    \item We use $\bar{M}_{\text{Pl}}$ only in ratios with other EDC parameters
    \item The relation $M_5^3 L = \bar{M}_{\text{Pl}}^2$ is a \textbf{structural} relation,
    not a numerical input
    \item We do NOT substitute the numerical value $\bar{M}_{\text{Pl}} = 2.435 \times 10^{18}$ GeV
    \item The final $G_F$ expression depends on $L$, $\sigma$, $\lambda$, not on $\bar{M}_{\text{Pl}}$ numerically
\end{enumerate}
\tagD
\end{proposition}

\subsection{Explicit Check}

Track B gives:
\begin{equation}
\label{eq:track-b-no-Mpl}
G_F = \text{function of } (L, \lambda, c_B, \text{spectral data})
\end{equation}

No $\bar{M}_{\text{Pl}}$ appears! \tagD

Track A gives:
\begin{equation}
\label{eq:track-a-with-Mpl}
g_4^2 = \frac{c_A}{L^{2/3} \bar{M}_{\text{Pl}}^{2/3}}
\end{equation}

But this can be rewritten using $M_5^3 L = \bar{M}_{\text{Pl}}^2$:
\begin{equation}
\label{eq:track-a-no-Mpl}
g_4^2 = \frac{c_A}{M_5 L}
\end{equation}

which depends on $M_5$ and $L$, not $\bar{M}_{\text{Pl}}$ directly. \tagD

%=============================================================================
\section{Dimensional and Convention Audit}
%=============================================================================

\subsection{Symbol Table}

\begin{table}[h]
\centering
\begin{tabular}{llll}
\toprule
Symbol & Dimension & Definition & Tag \\
\midrule
$g_5$ & $M^{-1/2}$ & 5D gauge coupling & \tagP \\
$g_4$ & $M^0$ & 4D gauge coupling & \tagD \\
$L$ & $M^{-1}$ & Compactification length & \tagP \\
$M_5$ & $M^1$ & 5D fundamental scale & \tagBL \\
$\bar{M}_{\text{Pl}}$ & $M^1$ & Reduced Planck mass & \tagBL \\
$\sigma$ & $M^4$ & Brane tension & \tagP \\
$\kappa_5$ & $M^{-3/2}$ & 5D gravitational coupling & \tagBL \\
$\lambda$ & $M^0$ & Pinning parameter & \tagDc/P \\
$k$ & $M^0$ ($\mathbb{Z}$) & Topological level & \tagP \\
$\beta$ & $M^0$ & EDC control parameter & \tagD \\
$b$ & $M^0$ & $\lambda\beta$ & \tagD \\
$m_n$ & $M^1$ & KK mass for mode $n$ & \tagD \\
$x_n$ & $M^0$ & $m_n L$ (dimensionless root) & \tagD \\
$G_F$ & $M^{-2}$ & Fermi constant & \tagD (target) \\
$c_A, c_B, c_C$ & $M^0$ & Track coefficients & \tagDc \\
$r_0, r_L$ & $M^{-1}$ & Brane kinetic lengths & \tagP \\
$\mathcal{I}_n$ & $M^0$ & Overlap integral & \tagD \\
\bottomrule
\end{tabular}
\caption{Symbol table with dimensions and tags.}
\label{tab:symbols}
\end{table}

\subsection{Dimensional Consistency Checks}

\begin{enumerate}
    \item $[g_5^2] = M^{-1}$: verified in Prop.~\ref{prop:dim-g5}
    \item $[g_4^2] = [g_5^2/L] = M^{-1}/M^{-1} = M^0$: dimensionless $\checkmark$
    \item $[G_F] = [g_4^2/m^2] = M^0/M^2 = M^{-2}$: $\checkmark$
    \item $[\beta] = [\sigma L^2/\bar{M}_{\text{Pl}}^2] = M^4 \cdot M^{-2}/M^2 = M^0$: $\checkmark$
\end{enumerate}

\tagI

%=============================================================================
\section{Open Items for Numerical Closure}
%=============================================================================

\subsection{Track-Independent Requirements}

\begin{enumerate}
    \item \textbf{Spectral roots $x_n(b)$}: Require numerical solution of BC equations
    \item \textbf{Overlap integrals $\mathcal{I}_n$}: Depend on fermion localization profile
    \item \textbf{Sum convergence}: Must verify $\sum_n |\mathcal{I}_n|^2/x_n^2 < \infty$
\end{enumerate}

\subsection{Track A Requirements}

\begin{enumerate}
    \item Coefficient $c_A$: undetermined at this level
    \item $M_5$ or equivalently $\sigma/\kappa_5^2$: requires additional input
\end{enumerate}

\subsection{Track B Requirements}

\begin{enumerate}
    \item Coefficient $c_B$: undetermined
    \item Selection of $k$: requires principle (from O-2/v37)
    \item $L$ value: requires closure (from v30 framework)
\end{enumerate}

\subsection{Track C Requirements}

\begin{enumerate}
    \item Coefficient $c_C$: constrained to $c_C \leq 1$ but not fixed
    \item Choice of $\Lambda_5$: $M_5$ or $\pi/L$ or other
\end{enumerate}

\subsection{Summary: What's Needed}

\begin{openbox}
\textbf{For $G_F$ numerical value}:
\begin{itemize}
    \item One of: $c_A$, $c_B$, or $c_C$ (track coefficient)
    \item Either $L$ or $M_5$ or equivalent scale
    \item Spectral data $x_n(b)$ from BC solver
    \item Overlap integrals from fermion profile
    \item Truncation: how many modes to sum
\end{itemize}
\tagOPEN
\end{openbox}

%=============================================================================
\section{Reviewer Trap Checklist}
%=============================================================================

\begin{table}[h]
\centering
\begin{tabular}{clcc}
\toprule
\# & Trap & Status & Resolution \\
\midrule
1 & Hidden $G_N$ via $\bar{M}_{\text{Pl}}$ & PASS & Sec.~10 trap check \\
2 & Hidden $\ell_P$ & PASS & Not used numerically \\
3 & $g_5$ dimension wrong & PASS & Prop.~\ref{prop:dim-g5} \\
4 & $g_4$ not dimensionless & PASS & Verified \\
5 & $\pi$-map inconsistency & PASS & Thm.~\ref{thm:pi-map} \\
6 & Brane terms ignored & PASS & Sec.~7 \\
7 & Forbidden inputs & PASS & None used \\
8 & Track A/B/C inconsistent & PASS & All give $[G_F] = M^{-2}$ \\
9 & Missing overlap integral & PASS & Included in bridge \\
10 & Normalization drift & PASS & Convention fixed \\
11 & Cutoff ambiguity & PASS & Two choices given \\
12 & $\lambda$ dimension wrong & PASS & Dimensionless \\
13 & Missing factor of 8 & PASS & From v34 \\
14 & KK mass formula wrong & PASS & Standard result \\
\bottomrule
\end{tabular}
\caption{Reviewer trap checklist.}
\label{tab:traps}
\end{table}

%=============================================================================
\section{Conclusions}
%=============================================================================

\subsection{What Was Derived}

\begin{enumerate}
    \item Dimensional structure: $[g_5^2] = M^{-1}$
    \item Track A: $g_5^2 = c_A/M_5$ from stiffness scaling
    \item Track B: $g_5^2 = 2\pi c_B L/\lambda$ from topological level
    \item Track C: $g_5^2 = 4\pi c_C/\Lambda_5$ from self-consistency
    \item Bridge formulas connecting each track to $G_F$
    \item $\pi$-map invariance theorem
    \item No hidden Planck trap verification
\end{enumerate}

\subsection{What Remains Open}

\begin{enumerate}
    \item Coefficient determination: $c_A$, $c_B$, $c_C$
    \item Scale input: $L$, $M_5$, or equivalent
    \item Spectral solver: $x_n(b)$ for given BC
    \item Overlap computation: $\mathcal{I}_n$ for given fermion profile
    \item Track selection: which mechanism is realized in nature
\end{enumerate}

%=============================================================================
\appendix
%=============================================================================

\section{Detailed Dimensional Analysis}
\label{app:dim}

\subsection{5D Action Dimensions}

Starting from:
\begin{equation}
\label{eq:app-action}
S = -\frac{1}{4g_5^2} \int d^5x \sqrt{-g}\, F_{MN}^2
\end{equation}

In natural units ($\hbar = c = 1$):
\begin{align}
\label{eq:app-d5x}
[d^5x] &= [d^4x][dy] = L^4 \cdot L = L^5 = M^{-5} \\
\label{eq:app-F}
[F_{MN}] &= [\partial A] = M \cdot M^{3/2} = M^{5/2} \\
\label{eq:app-F2}
[F_{MN}^2] &= M^5 \\
\label{eq:app-sqrt-g}
[\sqrt{-g}] &= 1 \quad \text{(flat metric)}
\end{align}

Wait, let me redo this more carefully. \tagI

\subsection{Gauge Field Dimension in 5D}

From canonical kinetic term, the gauge field must have:
\begin{equation}
\label{eq:app-A-dim}
[A_M] = M^{3/2}
\end{equation}

so that:
\begin{equation}
\label{eq:app-F-correct}
[F_{MN}] = [\partial_M A_N] = M \cdot M^{3/2} = M^{5/2}
\end{equation}

Therefore:
\begin{equation}
\label{eq:app-F2-correct}
[F_{MN}^2] = M^5
\end{equation}

From $[S] = M^0$:
\begin{equation}
\label{eq:app-g5-dim}
[g_5^{-2}] \cdot M^{-5} \cdot M^5 = M^0 \Rightarrow [g_5^{-2}] = M^0???
\end{equation}

This gives $g_5$ dimensionless, contradicting earlier! The issue: different
conventions for 5D vs 4D field normalization.

\subsection{Correct Convention}

The standard convention is:
\begin{equation}
\label{eq:app-convention}
[A_\mu] = M^1 \quad \text{in 4D}, \quad [A_M] = M^{3/2} \quad \text{in 5D}
\end{equation}

With this:
\begin{equation}
\label{eq:app-F-4d}
[F_{\mu\nu}] = M^2 \quad \text{(4D)}, \quad [F_{MN}] = M^{5/2} \quad \text{(5D)}
\end{equation}

But the action form \eqref{eq:5d-gauge-action} uses a different normalization where:
\begin{equation}
\label{eq:app-alt-conv}
[\tilde{A}_M] = M^1, \quad [\tilde{F}_{MN}] = M^2
\end{equation}

Then:
\begin{equation}
\label{eq:app-g5-final}
[g_5^{-2}] \cdot M^{-5} \cdot M^4 = M^0 \Rightarrow [g_5^{-2}] = M^1 \Rightarrow [g_5^2] = M^{-1}
\end{equation}

Confirmed! \tagI

\section{Chern-Simons Details}
\label{app:cs}

\subsection{5D Chern-Simons Form}

The 5D CS action is:
\begin{equation}
\label{eq:app-cs}
S_{\text{CS}} = \frac{k}{24\pi^2} \int_{\mathcal{M}_5} \omega_5^{\text{CS}}
\end{equation}

where $\omega_5^{\text{CS}}$ is the Chern-Simons 5-form:
\begin{equation}
\label{eq:app-omega5}
\omega_5 = \Tr\left(A \wedge dA \wedge dA - \frac{3i}{2} A \wedge A \wedge A \wedge dA
- \frac{3}{5} A \wedge A \wedge A \wedge A \wedge A\right)
\end{equation}

Under gauge transformation $A \to A + d\chi$:
\begin{equation}
\label{eq:app-gauge-var}
\delta S_{\text{CS}} = \frac{k}{24\pi^2} \int_{\partial\mathcal{M}_5} d\chi \wedge \omega_3
\end{equation}

where $\omega_3$ is related to the anomaly polynomial. \tagBL

\subsection{Level Quantization}

For consistency on manifolds with non-trivial $\pi_4(G)$:
\begin{equation}
\label{eq:app-level}
k \in \mathbb{Z}
\end{equation}

This is the origin of the topological level. \tagBL

\section{Self-Consistency Derivation}
\label{app:self}

\subsection{One-Loop Correction}

The gauge field self-energy at one loop:
\begin{equation}
\label{eq:app-self-energy}
\Pi_{\mu\nu}(p) = g_5^2 \int \frac{d^5k}{(2\pi)^5} V_{\mu\rho\sigma}(p,k,-k-p) \frac{1}{k^2} \frac{1}{(k+p)^2} V^{\rho\sigma}{}_\nu(-p,-k,k+p)
\end{equation}

Power counting:
\begin{equation}
\label{eq:app-power}
[\Pi] \sim g_5^2 \cdot M^5 \cdot M^{-4} = g_5^2 M
\end{equation}

The divergent part:
\begin{equation}
\label{eq:app-div}
\Pi_{\mu\nu}^{\text{div}} \sim g_5^2 \Lambda_5 (p^2 g_{\mu\nu} - p_\mu p_\nu)
\end{equation}

For perturbativity: $g_5^2 \Lambda_5 < 4\pi$. \tagBL

\section{Overlap Integral Computation}
\label{app:overlap}

\subsection{General Form}

The overlap integral:
\begin{equation}
\label{eq:app-overlap}
\mathcal{I}_n = \int_0^L dy\, w(y) |\chi_0(y)|^2 f_n(y)
\end{equation}

where:
\begin{itemize}
    \item $w(y)$ = warp factor (= 1 for flat)
    \item $\chi_0(y)$ = fermion zero-mode profile
    \item $f_n(y)$ = gauge KK mode profile
\end{itemize}

\subsection{Flat Space, Localized Fermion}

For $\chi_0 \propto e^{\alpha y}$ (localized) and $f_n = \sqrt{2/L}\cos(n\pi y/L)$:
\begin{equation}
\label{eq:app-overlap-loc}
\mathcal{I}_n = N^2 \sqrt{\frac{2}{L}} \int_0^L dy\, e^{2\alpha y} \cos\left(\frac{n\pi y}{L}\right)
\end{equation}

Using:
\begin{equation}
\label{eq:app-integral}
\int_0^L e^{ay} \cos(by) dy = \frac{e^{aL}(a\cos(bL) + b\sin(bL)) - a}{a^2 + b^2}
\end{equation}

We get a non-trivial $\mathcal{I}_n$ that alternates in sign. \tagD

\section{KK Spectrum Formulas}
\label{app:kk}

\subsection{Neumann/Neumann}

\begin{equation}
\label{eq:app-NN}
f_n(y) = \begin{cases}
1/\sqrt{L} & n = 0 \\
\sqrt{2/L}\cos(n\pi y/L) & n \geq 1
\end{cases}, \quad m_n = \frac{n\pi}{L}
\end{equation}

\subsection{Dirichlet/Dirichlet}

\begin{equation}
\label{eq:app-DD}
f_n(y) = \sqrt{\frac{2}{L}} \sin\left(\frac{n\pi y}{L}\right), \quad m_n = \frac{n\pi}{L}, \quad n \geq 1
\end{equation}

No zero-mode. \tagD

\subsection{Robin/Robin}

General boundary condition:
\begin{equation}
\label{eq:app-robin}
(\partial_y - m_b) f\big|_{y=0} = 0, \quad (\partial_y + m_b) f\big|_{y=L} = 0
\end{equation}

Spectrum determined by:
\begin{equation}
\label{eq:app-robin-spectrum}
\tan(m L) = \frac{2 m m_b}{m^2 - m_b^2}
\end{equation}

\tagD

\section{Epistemic Ledger}
\label{app:ledger}

\begin{table}[h]
\centering
\begin{tabular}{llc}
\toprule
Item & Tag & Eq./Prop \\
\midrule
5D gauge action & \tagP & \eqref{eq:5d-gauge-action} \\
$[g_5^2] = M^{-1}$ & \tagD & Prop.~\ref{prop:dim-g5} \\
$g_4^2 = g_5^2/L$ (flat) & \tagD & \eqref{eq:g4-g5-simple} \\
Track A: $g_5^2 = c_A/M_5$ & \tagDc & Prop.~\ref{prop:track-a-formula} \\
Track B: $g_5^2 = 2\pi c_B L/\lambda$ & \tagDc/P & Prop.~\ref{prop:track-b} \\
Track C: $g_5^2 = 4\pi c_C/\Lambda_5$ & \tagP$\to$Dc & Prop.~\ref{prop:track-c} \\
Bridge formulas & \tagD & Sec.~8 \\
$\pi$-map invariance & \tagD & Thm.~\ref{thm:pi-map} \\
No Planck trap & \tagD & Prop.~\ref{prop:no-trap} \\
$M_5^3 L = \bar{M}_{\text{Pl}}^2$ & \tagBL & \eqref{eq:M5-Mpl} \\
$\lambda = |k|/(2\pi)$ & \tagDc/P & \eqref{eq:lambda-k} \\
Track coefficients & \tagOPEN & $c_A, c_B, c_C$ \\
\bottomrule
\end{tabular}
\caption{Epistemic ledger for v36.}
\label{tab:ledger}
\end{table}

\section{Extended Track Analysis}
\label{app:extended}

\subsection{Track A: Alternative Derivation}

Starting from the Einstein-Hilbert action in 5D:
\begin{equation}
\label{eq:app-eh-5d}
S_{\text{grav}} = \frac{1}{2\kappa_5^2} \int d^5x \sqrt{-g} \mathcal{R}
\end{equation}

where $[\kappa_5^2] = M^{-3}$. \tagBL

The relation to $M_5$:
\begin{equation}
\label{eq:app-kappa-M5}
\kappa_5^2 = \frac{1}{M_5^3}
\end{equation}

Combined with brane tension $\sigma$:
\begin{equation}
\label{eq:app-M5-sigma}
M_5^3 = \sigma \cdot \kappa_5^2 \cdot \text{(dimensionless)}
\end{equation}

This gives:
\begin{equation}
\label{eq:app-M5-explicit}
M_5 = (\sigma \kappa_5^2)^{1/3}
\end{equation}

\subsection{Track A: Coupling as Stiffness Ratio}

Define the stiffness ratio:
\begin{equation}
\label{eq:app-stiff-ratio}
\mathcal{R}_{\text{stiff}} = \frac{\sigma}{M_5^4} = \sigma^{1/3} \kappa_5^{8/3}
\end{equation}

The gauge coupling can be written:
\begin{equation}
\label{eq:app-g5-stiff}
g_5^2 = \frac{c_A}{M_5} = c_A \cdot \frac{\kappa_5^{2/3}}{\sigma^{1/3}}
\end{equation}

In terms of the stiffness ratio:
\begin{equation}
\label{eq:app-g5-R}
g_5^2 = c_A \cdot M_5^{-1} = c_A \cdot \mathcal{R}_{\text{stiff}}^{1/4} \cdot \sigma^{-1/4}
\end{equation}

\tagD

\subsection{Track B: Discrete vs Continuous}

For continuous $k$ (non-quantized):
\begin{equation}
\label{eq:app-cont-k}
g_5^2 = \frac{c_B' L}{k}
\end{equation}

The 4D coupling:
\begin{equation}
\label{eq:app-g4-cont}
g_4^2 = \frac{c_B'}{k}
\end{equation}

This allows continuously variable coupling, unlike the quantized case. \tagP

\subsection{Track B: Level from Anomaly Matching}

If the anomaly polynomial requires specific level:
\begin{equation}
\label{eq:app-anomaly-k}
k = k_{\text{anomaly}} = n_f - n_b
\end{equation}
where $n_f$ and $n_b$ count chiral fermion and boson contributions. \tagBL

This fixes $k$ from matter content. \tagD

\subsection{Track C: Two-Loop Extension}

At two loops, the self-consistency bound becomes:
\begin{equation}
\label{eq:app-2loop}
g_5^2 \Lambda_5 + (g_5^2 \Lambda_5)^2 / (16\pi^2) < 4\pi
\end{equation}

Solving for $g_5^2$:
\begin{equation}
\label{eq:app-g5-2loop}
g_5^2 < \frac{4\pi}{\Lambda_5} \cdot \frac{1}{1 + \Lambda_5/(4\pi)}
\end{equation}

For $\Lambda_5 \ll 4\pi$, this reduces to the one-loop result. \tagD

\section{Warped Space Generalization}
\label{app:warped}

\subsection{RS-like Metric}

For warped metric:
\begin{equation}
\label{eq:app-rs-metric}
ds^2 = e^{-2\sigma(y)} \eta_{\mu\nu} dx^\mu dx^\nu + dy^2
\end{equation}

with $\sigma(y) = k|y|$ (RS1). \tagBL

\subsection{Gauge Action in Warped Space}

\begin{equation}
\label{eq:app-warped-action}
S_{\text{gauge}}^{\text{warped}} = -\frac{1}{4g_5^2} \int d^4x \int_0^L dy\, e^{-2\sigma} F_{\mu\nu}^2
+ e^{-4\sigma} F_{\mu 5}^2
\end{equation}

The different warp factors for different components affect the 4D coupling:
\begin{equation}
\label{eq:app-g4-warped}
\frac{1}{g_4^2} = \frac{1}{g_5^2} \int_0^L dy\, e^{-2\sigma(y)} |f_0(y)|^2
\end{equation}

\tagD

\subsection{Zero-Mode Profile}

In RS1:
\begin{equation}
\label{eq:app-f0-rs}
f_0(y) = \frac{e^{kL} - 1}{kL(e^{2kL} - 1)^{1/2}} e^{ky}
\end{equation}

Localized toward UV brane. \tagBL

\subsection{4D Coupling in RS1}

\begin{equation}
\label{eq:app-g4-rs}
g_4^2 = g_5^2 \cdot \frac{2k}{e^{2kL} - 1}
\end{equation}

For $kL \gg 1$:
\begin{equation}
\label{eq:app-g4-rs-large}
g_4^2 \approx 2k g_5^2 e^{-2kL}
\end{equation}

Exponentially suppressed! \tagD

\subsection{Implications for Tracks}

\textbf{Track A in RS1:}
\begin{equation}
\label{eq:app-track-a-rs}
g_4^2 = \frac{2k c_A}{M_5} e^{-2kL}
\end{equation}

\textbf{Track B in RS1:}
\begin{equation}
\label{eq:app-track-b-rs}
g_4^2 = \frac{4\pi k c_B}{\lambda} e^{-2kL}
\end{equation}

The exponential factor provides additional suppression. \tagD

\section{Connection to Gauge Unification}
\label{app:unification}

\subsection{GUT-Scale Coupling}

At the GUT scale $M_{\text{GUT}}$:
\begin{equation}
\label{eq:app-gut-coupling}
g_{\text{GUT}}^2 = g_1^2(M_{\text{GUT}}) = g_2^2(M_{\text{GUT}}) = g_3^2(M_{\text{GUT}})
\end{equation}

In 5D unification, this becomes:
\begin{equation}
\label{eq:app-5d-unif}
g_{\text{GUT}}^2 = \frac{g_5^2}{L}
\end{equation}

if compactification scale $\sim M_{\text{GUT}}$. \tagBL

\subsection{Track Implications}

\textbf{Track B prediction:}
\begin{equation}
\label{eq:app-track-b-gut}
g_{\text{GUT}}^2 = \frac{2\pi c_B}{\lambda} = \frac{(2\pi)^2 c_B}{|k|}
\end{equation}

For $g_{\text{GUT}}^2 \approx 0.5$ (typical GUT value):
\begin{equation}
\label{eq:app-k-gut}
|k| \approx \frac{(2\pi)^2 c_B}{0.5} \approx 80 c_B
\end{equation}

If $c_B \sim 1$, then $k \sim 80$. \tagD

\subsection{Running Below Compactification}

Below $M_c = \pi/L$, the coupling runs as 4D:
\begin{equation}
\label{eq:app-running}
\frac{1}{g_4^2(\mu)} = \frac{1}{g_4^2(M_c)} + \frac{b}{8\pi^2} \ln\frac{M_c}{\mu}
\end{equation}

where $b$ is the beta function coefficient. \tagBL

\section{Numerical Estimates}
\label{app:numeric}

\subsection{Track B Example}

With parameters:
\begin{align}
\label{eq:app-params}
L &= 10^{-17} \text{ m (symbolic)} \\
\lambda &= 1 \quad (k = 2\pi) \\
c_B &= 1
\end{align}

The 4D coupling:
\begin{equation}
\label{eq:app-g4-example}
g_4^2 = 2\pi \cdot 1 / 1 = 2\pi \approx 6.28
\end{equation}

This is $O(1)$, consistent with gauge couplings. \tagD

\subsection{$G_F$ Order of Magnitude (Track B)}

Using $m_1 = \pi/L$ and $|\mathcal{I}_1|^2 \sim 1$:
\begin{equation}
\label{eq:app-GF-oom}
G_F \sim \sqrt{2} \cdot \frac{2\pi c_B}{\lambda} \cdot \frac{L^2}{\pi^2} = \frac{2\sqrt{2} c_B L^2}{\pi \lambda}
\end{equation}

With $[G_F] = M^{-2}$ and $[L^2] = M^{-2}$: consistent! \tagI

\subsection{Scaling Test}

For $L \to 2L$:
\begin{equation}
\label{eq:app-scale-test}
G_F \to 4 G_F
\end{equation}

$G_F \propto L^2$: larger extra dimension gives larger Fermi constant. \tagD

\section{Alternative Formulations}
\label{app:alt}

\subsection{6D Compactification}

In 6D on $T^2$ with radii $R_5, R_6$:
\begin{equation}
\label{eq:app-6d-action}
S_6 = -\frac{1}{4g_6^2} \int d^6x \sqrt{-g}\, F_{MN}^2
\end{equation}

where $[g_6^2] = M^{-2}$. \tagP

The 4D coupling:
\begin{equation}
\label{eq:app-g4-6d}
g_4^2 = \frac{g_6^2}{(2\pi)^2 R_5 R_6}
\end{equation}

\tagD

\subsection{Orbifold vs Interval}

On orbifold $S^1/\mathbb{Z}_2$:
\begin{equation}
\label{eq:app-orbifold}
g_4^2 = \frac{g_5^2}{\pi R}
\end{equation}

On interval $[0, L]$:
\begin{equation}
\label{eq:app-interval}
g_4^2 = \frac{g_5^2}{L}
\end{equation}

With $L = \pi R$, these agree. \tagDc

\subsection{Non-Abelian vs Abelian}

For non-Abelian gauge group with structure constants $f^{abc}$:
\begin{equation}
\label{eq:app-nonab}
[A_\mu, A_\nu] \to \frac{1}{g_5} [A_\mu, A_\nu]
\end{equation}

The coupling appears in commutator. Same dimensional analysis applies. \tagD

%=============================================================================
\end{document}
